\documentclass[11pt]{article}
\usepackage[margin=1in]{geometry}
\usepackage{amsmath,amssymb,amsthm,mathtools}
\usepackage{graphicx}
\usepackage{microtype}
\usepackage{xcolor}
\usepackage{enumitem}
\usepackage[colorlinks=true,linkcolor=blue!55!black,citecolor=blue!55!black,urlcolor=blue!55!black]{hyperref}
\setlength{\emergencystretch}{2em}

\newtheorem{theorem}{Theorem}
\newtheorem{lemma}{Lemma}
\newcommand{\R}{\mathbb R}
\newcommand{\F}{\mathcal F}
\newcommand{\Dminus}{{}^{-}\!\mathcal D^{\alpha}}
\newcommand{\Dplus}{{}^{+}\!\mathcal D^{\alpha}}
\newcommand{\Wminus}{{}^{-}\!W^{\alpha,p}(\R)}
\newcommand{\Wplus}{{}^{+}\!W^{\alpha,p}(\R)}
\newcommand{\Gag}{\widetilde W^{\alpha,p}(\R)}

\title{A counterexample to the conjectured Gagliardo equivalence\\
for one-sided weak fractional Sobolev spaces}
\author{Candidate result for arXiv:2004.10833}
\date{26 August 2026}

\begin{document}
\maketitle

\begin{center}
\fcolorbox{black!35}{blue!4}{\parbox{0.91\linewidth}{
\textbf{Status: candidate counterexample, likely valid.}
For every $0<\alpha<1$ and every $1<p<\infty$, $p\ne2$, the conjectured
equality on $\R$ is false.  The endpoint cases $p=1,\infty$ are not claimed.
}}
\end{center}

\section{Source and exact conjecture}

The source is X.~Feng and M.~Sutton, \emph{A new theory of fractional
differential calculus and fractional Sobolev spaces: One-dimensional case},
arXiv:2004.10833v2 \cite{FSlong}.  The conjecture appears in Remark 5.19(c),
printed page 64 (source line 3960).  It is repeated verbatim in the authors'
concise follow-up arXiv:2007.10245v1 \cite{FSshort}.

\begin{figure}[h]
  \centering
  \includegraphics[width=\textwidth]{figures/open_problem_crop.png}
  \caption{The source conjecture on printed page 64 of arXiv:2004.10833v2.
  This is a direct crop of the archived source PDF.}
\end{figure}

The relevant clause reads:
\begin{quote}
``We conjecture that ${}^{\pm}W^{\alpha,p}(\R)\ne
\widehat W^{\alpha,p}(\R)$, but ${}^{\pm}W^{\alpha,p}(\R)=
\widetilde W^{\alpha,p}(\R)$ for $p\ne2$ and $0<\alpha<1$.''
\end{quote}
We disprove the asserted equality with $\widetilde W^{\alpha,p}(\R)$ for
every interior exponent $p\ne2$.  We do not address the separate assertion
involving $\widehat W^{\alpha,p}$.

\section{Definitions and proof intuition}

For $0<\alpha<1$, the source defines $\Wminus$ and $\Wplus$ by requiring
$u\in L^p(\R)$ and, respectively, the left or right weak fractional derivative
to lie in $L^p(\R)$.  The Gagliardo space $\Gag$ consists of the functions
$u\in L^p(\R)$ for which
\[
 [u]_{\Gag}^p
 =\int_{\R}\!\int_{\R}\frac{|u(x)-u(y)|^p}
 {|x-y|^{1+\alpha p}}\,dx\,dy<\infty.
\]

\subsection*{Proof Intuition}

The two norms aggregate dyadic frequencies in different ways.  On the line,
the left and right fractional derivatives have symbols $(i\xi)^\alpha$ and
$(-i\xi)^\alpha$.  After removing a harmless Hilbert-transform phase, either
graph norm is the Bessel-potential norm $F^\alpha_{p,2}$, whose dyadic pieces
are summed in $\ell^2$ pointwise.  The Gagliardo norm is the Besov norm
$B^\alpha_{p,p}$, whose dyadic pieces are summed in $\ell^p$.  A sequence of
widely separated frequency packets with a common spatial envelope turns the
two memberships into the elementary alternatives $a\in\ell^2$ and
$a\in\ell^p$.  These agree only at $p=2$.

\section{The counterexample theorem}

\begin{theorem}\label{thm:main}
Let $0<\alpha<1$ and $1<p<\infty$.  Then, with equivalent norms,
\[
 \Wminus=\Wplus=H^{\alpha,p}(\R)=F^\alpha_{p,2}(\R),
 \qquad
 \Gag=B^\alpha_{p,p}(\R).
\]
Consequently, if $p\ne2$, neither one-sided space equals $\Gag$.
More precisely:
\begin{enumerate}[label=\textup{(\roman*)}]
\item if $2<p<\infty$, there exists a real
$u\in\Gag\setminus(\Wminus\cup\Wplus)$;
\item if $1<p<2$, there exists a real
$u\in(\Wminus\cap\Wplus)\setminus\Gag$.
\end{enumerate}
\end{theorem}

We split the proof into the operator identification and the packet
construction.

\begin{lemma}[The source graph spaces are Bessel-potential spaces]
\label{lem:graph}
For $0<\alpha<1$ and $1<p<\infty$, each one-sided weak derivative graph norm
is equivalent to
\[
 \|u\|_{L^p}+\|(-\Delta)^{\alpha/2}u\|_{L^p}.
\]
Thus each graph space is $H^{\alpha,p}(\R)=F^\alpha_{p,2}(\R)$.
\end{lemma}

\begin{proof}
Use the principal branches
\[
 (i\xi)^\alpha=|\xi|^\alpha e^{i\theta\operatorname{sgn}\xi},
 \qquad
 (-i\xi)^\alpha=|\xi|^\alpha e^{-i\theta\operatorname{sgn}\xi},
 \qquad \theta=\frac{\pi\alpha}{2}.
\]
The source defines the left weak derivative by duality against the right
classical derivative, and conversely.  Fourier duality on test functions gives
\[
 \F(\Dminus u)(\xi)=(i\xi)^\alpha\widehat u(\xi),
 \qquad
 \F(\Dplus u)(\xi)=(-i\xi)^\alpha\widehat u(\xi)
\]
in tempered distributions.  This identity is particularly unambiguous for
the $L^1\cap L^p$ witnesses constructed below; conversely an $L^p$ inverse
Fourier multiplier satisfying the identity is the source weak derivative by
the defining test relation.

Let $H$ be the Hilbert transform, with multiplier $-i\operatorname{sgn}\xi$.
The phase multipliers above are
\[
 e^{\pm i\theta\operatorname{sgn}\xi}
 =\cos\theta\mp\sin\theta\,\widehat H(\xi).
\]
They and their inverses are bounded on $L^p$ by the M.~Riesz theorem
\cite{Stein}.  It follows that
\[
 \|\Dminus u\|_p\asymp\|(-\Delta)^{\alpha/2}u\|_p
 \asymp\|\Dplus u\|_p.
\]
The inhomogeneous graph norm on the right is equivalent, by the usual smooth
multiplier comparison, to
$\|(1-\Delta)^{\alpha/2}u\|_p$, which defines $H^{\alpha,p}$.  The classical
Littlewood--Paley characterization gives
$H^{\alpha,p}=F^\alpha_{p,2}$ \cite{Triebel}.
\end{proof}

\begin{lemma}[A sparse packet separates the sequence exponents]
\label{lem:packets}
Fix $0\ne\eta\in\mathcal S(\R)$ whose Fourier transform is supported in a
sufficiently small interval about the origin.  Put $n_k=4k$.  For every
polynomially bounded scalar sequence $a=(a_k)$, the series
\[
 u_a(x)=\sum_{k=1}^\infty a_k,2^{-\alpha n_k}
 e^{i2^{n_k}x}\eta(x)                                      \tag{1}
\]
converges in $L^1(\R)\cap L^p(\R)$ and satisfies
\[
 u_a\in F^\alpha_{p,2}(\R)\iff a\in\ell^2,
 \qquad
 u_a\in B^\alpha_{p,p}(\R)\iff a\in\ell^p.               \tag{2}
\]
\end{lemma}

\begin{proof}
The factor $2^{-\alpha n_k}$ makes (1) absolutely summable in both $L^1$ and
$L^p$.  The Fourier supports of its terms lie in pairwise separated dyadic
annuli.  Choose smooth projections $Q_k$, uniformly obtained by dilation,
which equal one on those supports and vanish on the other packet supports.
Admissible Littlewood--Paley norms are independent of the chosen smooth
dyadic resolution \cite{Triebel}.  Hence the sparse square-function and Besov
norms are, with constants independent of $a$,
\begin{align*}
 \|u_a\|_{F^\alpha_{p,2}}
 &\asymp \|u_a\|_p+
 \left\|\left(\sum_k
 |2^{\alpha n_k}Q_ku_a|^2\right)^{1/2}\right\|_p\\
 &=\|u_a\|_p+\|\eta\|_p\left(\sum_k|a_k|^2\right)^{1/2},\\
 \|u_a\|_{B^\alpha_{p,p}}
 &\asymp \|u_a\|_p+
 \left(\sum_k2^{\alpha p n_k}\|Q_ku_a\|_p^p\right)^{1/p}\\
 &=\|u_a\|_p+\|\eta\|_p\left(\sum_k|a_k|^p\right)^{1/p}.
\end{align*}
This proves (2).  Equivalently, one may insert the $Q_k$ into any fixed
dyadic resolution; finite overlap and the vector-valued multiplier theorem
give the same two estimates.
\end{proof}

\begin{proof}[Proof of Theorem \ref{thm:main}]
The first identity is Lemma \ref{lem:graph}.  The classical difference
characterization of Besov spaces gives
$\Gag=B^\alpha_{p,p}(\R)$ for $0<\alpha<1$ and
$1<p<\infty$ \cite{Triebel}.

Set
\[
 \beta=\frac12\left(\frac12+\frac1p\right),
 \qquad a_k=k^{-\beta}.
\]
If $p>2$, then $1/p<\beta<1/2$, so
$a\in\ell^p\setminus\ell^2$.  Lemma \ref{lem:packets} gives
$u_a\in B^\alpha_{p,p}\setminus F^\alpha_{p,2}$.  If $p<2$, then
$1/2<\beta<1/p$, so $a\in\ell^2\setminus\ell^p$, and the reverse membership
holds.

The displayed witness is complex-valued only for notational convenience.
Both spaces are stable under real and imaginary parts.  If both components
belonged to the excluded space, then their complex sum would also belong to
it.  Thus at least one of $\operatorname{Re}u_a$ and
$\operatorname{Im}u_a$ is a real-valued witness with the stated membership.
\end{proof}

\section{Verification, limitations, and novelty}

The proof uses four classical tools: the $L^p$ bound for the Hilbert
transform; the identity $H^{\alpha,p}=F^\alpha_{p,2}$; the difference
description $\widetilde W^{\alpha,p}=B^\alpha_{p,p}$; and invariance under
the choice of smooth dyadic resolution.  These are all standard for interior
$p$ \cite{Stein,Triebel}.  No finite computation can verify the separation:
every finite packet sum lies in both spaces, and the obstruction is exactly
the infinite $\ell^2$ versus $\ell^p$ summability.

The endpoint cases $p=1$ and $p=\infty$ remain outside the theorem because
the Hilbert transform is not bounded on those $L^p$ spaces.  We also make no
claim about the source's nonstandard Fourier-$L^p$ space
$\widehat W^{\alpha,p}$, nor about finite intervals in Remark 5.19(d).

A bounded search (26 August 2026) covered the run's registry, solution,
attempt, and proof-gap indexes; the exact wording of Remark 5.19(c); the
repeated arXiv:2007.10245 wording; exact-title, author, and core-term arXiv
queries; and the later visible Feng--Sutton arXiv work.  No later paper was
found that explicitly claims to resolve this exact conjecture.  Novelty
confidence is therefore moderate, not high: the general distinction between
$F^\alpha_{p,2}$ and $B^\alpha_{p,p}$ is classical.  The candidate
contribution is the exact identification of the source graph spaces and the
explicit application that disproves its conjecture for all interior
$p\ne2$.

\paragraph{Human review recommendation.}
Send to a harmonic analyst or function-space expert.  The most useful checks
are the Fourier-duality convention in Lemma \ref{lem:graph} and the sparse
projection norm equivalence in Lemma \ref{lem:packets}.

\begin{thebibliography}{9}
\bibitem{FSlong}
X.~Feng and M.~Sutton,
\emph{A new theory of fractional differential calculus and fractional Sobolev
spaces: One-dimensional case}, arXiv:2004.10833v2 (2020).

\bibitem{FSshort}
X.~Feng and M.~Sutton,
\emph{On New Families of Fractional Sobolev Spaces},
arXiv:2007.10245v1 (2020).

\bibitem{Stein}
E.~M.~Stein,
\emph{Singular Integrals and Differentiability Properties of Functions},
Princeton University Press, 1970.

\bibitem{Triebel}
H.~Triebel,
\emph{Theory of Function Spaces},
Birkh\"auser, 1983.
\end{thebibliography}

\end{document}
